\chapter{The Limits of Next-Token Prediction}

\section{Granting Everything and Still Falling Short}

Suppose, for the sake of argument, that Chapter 3's problem is entirely solved. Suppose the system generating the courtyard has perfect memory: every fact, once established, is guaranteed to be written; every relevant fact, once written, is guaranteed to be retrieved when needed; every contradiction is guaranteed to be reconciled the moment it would otherwise arise. This is a considerable concession, well beyond what any deployed retrieval system currently offers, and it is worth granting it fully before asking the question this chapter exists to ask. Given all of this, given perfect access to everything that has ever been established about the wall, does the system that generates the wall's next appearance actually produce a persistent world?

The honest answer is not automatically. Perfect information, available to a generative process, does not by itself guarantee that the process uses that information as a binding constraint rather than as one more influence, however strong, on a fundamentally different kind of computation. This chapter argues that the generative principle underlying nearly every system this book has discussed, next-token prediction and its close relatives across other modalities, is built to answer a different question than the one persistence actually requires, and that this mismatch survives even a perfect solution to the storage problem Chapters 2 and 3 addressed.

\section{Two Kinds of "What Comes Next"}

It is worth distinguishing two computations that are easy to conflate because they frequently produce the same answer. Statistical continuation asks: given everything currently in context, what is the most probable next token, frame, or region of an image, according to a distribution learned from training data. Causal continuation asks a different question: given the actual, specific prior state of a system, what must or should follow, according to whatever rules actually govern that system's evolution. A physics engine computes causal continuation: given a ball's exact position and velocity, its next position is determined, not merely likely, by the equations governing motion. A video model trained on footage of falling balls computes statistical continuation: given the frames so far, what continuation looks like a plausible falling ball, according to patterns extracted from many other falling balls it was trained on. These are different computations, arrived at through entirely different means, and it is only a contingent, empirical fact when they happen to agree.

\section{Why the Two Usually Agree, and Why the Agreement Is Misleading}

In practice, statistical and causal continuation agree often enough that the distinction can seem like a philosopher's nicety rather than a live engineering concern. A model trained on enough footage of falling balls learns a distribution that closely tracks actual physics, not because it has represented Newton's laws anywhere internally, but because physically accurate continuations are, unsurprisingly, the statistically common ones in a large enough corpus of real footage. This agreement is real, and it is also entirely contingent on training data densely covering the situation actually being generated. It breaks down precisely where this book has been focused from Chapter 1 onward: at the specific, arbitrary, already-established facts that a persistent world needs to honor, and that no amount of generic plausibility can substitute for. The wall's crack sits in a specific location because it was established there, once, by a specific prior act of invention, not because that specific location is any more statistically probable than any of the countless other locations a crack could plausibly appear. Statistical continuation has no privileged reason to reproduce that specific location. Only causal continuation, actually reading off the established fact rather than sampling from a general distribution over plausible cracks, gets it right reliably rather than by coincidence.

\section{Retrieved Memory Does Not Make Prediction Causal}

This is where Chapter 3's concession, granted in full at the start of this chapter, turns out not to be enough, and the reason is worth stating precisely because it is easy to miss. Even with the crack's exact position perfectly retrieved and placed directly into the model's context, an autoregressive predictor still generates its output by asking what is statistically likely to come next given this context, not what is required to come next given this context. Having the fact present in context makes the correct continuation highly probable, since a well-trained model will generally learn to weight in-context evidence heavily. Highly probable is not the same as required. The retrieved fact is one strong influence among many competing influences drawn from the model's training distribution, and nothing about the underlying computation prevents that distribution's other pressures, stylistic tendencies, competing plausible details, the accumulated weight of countless other walls the model was trained on, from occasionally winning out, particularly under distribution shift, adversarial or unusual framing, or simply across enough generation steps that low-probability deviation eventually occurs somewhere. This is a different failure than anything Chapters 2 and 3 described. Those chapters were about whether the right information makes it into context at all. This chapter is about what happens once it is there: a fact present in context is honored with high probability, not with certainty, because the mechanism converting context into output was never built to enforce anything. It was built to continue plausibly, and plausible, however well-informed, is not the same claim as necessary.

\section{What Causal Continuation Would Actually Require}

It is worth sketching, without yet building it, what would need to be different for continuation to be causal rather than statistical. The next state of a persistent world cannot merely be conditioned on the current state, weighted alongside everything else a generative process has learned to find plausible. It needs to be constrained by the current state, in a sense a probability distribution over continuations does not, by itself, provide: some notion of which continuations remain admissible given what has actually been established, and a mechanism capable of ruling out inadmissible continuations outright rather than merely assigning them somewhat lower probability. A system that only makes the wrong crack-location less likely has not solved this chapter's problem. A system that makes it impossible has. This book will spend a considerable portion of its remaining chapters building exactly this distinction into precise, workable form; for now, it is enough to have located the gap precisely: not in what a generative system knows, which Chapter 3's concession granted could in principle be perfect, but in what a generative system is architecturally capable of enforcing once it knows it.

\section{Looking Ahead}

The last three chapters have each isolated one piece of why current systems fail to sustain a persistent world: buffers that cannot distinguish forgetting from never-happening, document stores that cannot reconcile their own contents, and a generative principle that treats even perfectly retrieved facts as merely probable rather than binding. Chapter 5 draws these threads together by looking, in detail, at systems where this entire chain of failures is visible end to end, including one built simply enough that its mechanism can be inspected directly rather than inferred, to show concretely what generation without an underlying reality actually looks like in practice.
